\documentclass{article}
\usepackage[left=2cm, right=2cm, top=2cm, bottom=2cm]{geometry}
\usepackage{graphicx}
\usepackage{amsmath}
\usepackage{amssymb}
\usepackage{amsthm}
\usepackage{fancyhdr}
\usepackage{verbatim}
\usepackage{listings}
\usepackage{xcolor}
\usepackage{pdfpages}
\usepackage{cancel}
\usepackage{physics}

\lstset{
    language=C++,
    basicstyle=\ttfamily\footnotesize,
    keywordstyle=\color{blue},
    commentstyle=\color{green},
    stringstyle=\color{red},
    numbers=left,
    numberstyle=\tiny\color{gray},
    frame=single,
    breaklines=true,
    captionpos=b,
}


\title{MTH 553: Lecture \# 16}
\author{Cliff Sun}

\newtheorem{theorem}{Theorem}[section]
\newtheorem{lemma}[theorem]{Lemma}
\newtheorem{definition}[theorem]{Definition}
\newtheorem{conjecture}[theorem]{Conjecture}
\newtheorem{proposition}[theorem]{Proposition}
\newtheorem{corollary}[theorem]{Corollary}
\newtheorem{one minute paper}[theorem]{One Minute Paper}

\pagestyle{fancy}
\lhead{\textbf{\thepage}\ \ \nouppercase{\rightmark}}
\chead{MTH 553: Lecture \# 16}
\rhead{Cliff Sun}

\begin{document}

\maketitle

\section*{Lecture Span}
\begin{enumerate}
    \item Energy method continued
\end{enumerate}

\section*{Lower order terms}

Take $c=1$. We consider \underline{Dispersion} (waves don't all travel at the same speed). 

\subsection*{Dispersion}

Add a zero order term to the HWE: 

\begin{equation}
    u_{tt} - u_{xx} + \lambda u = 0
\end{equation}

Where $\lambda \geq 0$ is a constant. This means that 

\begin{equation}
    u_{tt} = u_{xx} - \lambda u
\end{equation}

Where $\lambda u$ opposes the acceleration, i.e. a restoring force. Let $k \in \mathbb{R}$ and $\omega > 0$. Try $u(x,t) = e^{i(kx - \omega t)}$. 
Here, $k$ is the spatial frequency and $\omega$ is the temporal frequency. Plugging this into equation (2) and obtain 

\begin{equation}
    \omega = \sqrt{k^2 + \lambda} 
\end{equation}

This is called the dispersion relation, where the temporal frequency depends on the spatial frequency. Note that $u$ is a left 
or right moving wave with velocity $$c(k) = \omega/k = \frac{\sqrt{k^2 + \lambda}}{k}$$

\begin{enumerate}
    \item $\lambda = 0$, and $c(k) = 1$. Therefore, all waves travel at the same speed
    \item $\lambda > 0$, then $c(k)$ changes with $k$ and therefore not all waves travel at the same speed. Therefore, this PDE is called dispersive. 
\end{enumerate}

\subsection*{Example}

Assume $\lambda > 0$. Suppose initial condition, then use Fourier transform to write down an equation. Assume the following initial condition:

\begin{equation}
    u(x,0) = \int_{-\infty}^{\infty}A(k)e^{ikx}dk
\end{equation}

Solution, 

\begin{equation}
    u(x,t) = \int_{-\infty}^{\infty}A(k)e^{i(kx - \omega(k) t)}dk
\end{equation}

Observe, even if $u(x,0)$ has compact support. Then, the solution $u(x,t)$ will not since the low frequency components of the initial 
data travel arbitrarily fast. This is because $c(k) \rightarrow \infty$ as $k \rightarrow 0$. 

\subsection*{Schr\"odinger equation}

\begin{equation}
    u_t = iu_{xx}
\end{equation}

Where $u(x)$ is complex valued. 

\begin{proof}
    Let $u= e^{i(kx-\omega t)}$. And we obtain 
    \begin{equation}
        \omega = k^2
    \end{equation}
    And therefore, velocity is $\omega/k = k$. 
\end{proof}

\textbf{Remark:} Dispersive equations can still conserve quantities. For instance, equation (1) has the energy 

\begin{equation}
    \epsilon(t) = \frac{1}{2}\int_{-\infty}^{\infty}\left( u_t^2 + u_x^2 + \lambda u^2 \right)dx
\end{equation}

Check that $\epsilon'(t) = 0$ assuming boundary terms vanish (compact support). Schr\"odinger conserves the $L^2$ norm. In other words, 

\begin{equation}
    \frac{d}{dt}\left( \int_{\mathbb{R}}|u|^2 dx \right) = 0
\end{equation}

This represents the total probability of finding the particle.

\subsection*{Dissipation}

We call an equation \underline{dissipative} if the energy decreases with time. As an example, consider the equation

\begin{equation}
    u_{tt} + \alpha u_t - u_{xx} = 0
\end{equation}

This $\alpha u_t$ term is a friction term. How does this affect the energy? Take the usual energy:

\begin{equation}
    \epsilon(t) = \frac{1}{2}\int_{-\infty}^{\infty}\left( u_t^2 + u_x^2 \right)dx
\end{equation}

Then, 

\begin{align*}
    \epsilon'(t) &= \int \left( u_t u_{tt} + u_x u_{xt} \right) \\
    &= \int \nabla \cdot (u_t u_x) - \alpha \int u_t^2 < 0
\end{align*}

Such a result has the effect of controlling $u_t$ and $u_x$. We can also reduce the lower order terms in general, suppose 

\begin{equation}
    u_{tt} - c^u_{xx} + \alpha u_t + \beta u_x + \gamma u =0
\end{equation}

Then let 

\begin{equation}
    v(x,t) = e^{(\alpha t - \beta c^{-2}x)/2}u(x,t)
\end{equation}

Then the PDE transforms into 

\begin{equation}
    v_{tt} - c^2 v_{xx} + \left( -\frac{\alpha^2}{4} + \frac{\beta^2}{4c^2} + \gamma \right) v = 0
\end{equation}

\end{document}